\documentclass[11pt]{article}
\usepackage[a4paper,margin=25mm]{geometry}
\usepackage{amsmath,amssymb,amsthm,hyperref}
\newtheorem{theorem}{Theorem}
\newtheorem{lemma}{Lemma}
\newtheorem{corollary}{Corollary}
\title{Consecutive attraction for augmented happy functions}
\author{Principia Math In Progress}
\date{28 September 2026}
\begin{document}
\maketitle
\begin{abstract}
For every odd base $b>2$ and odd augmenting constant $c>0$, each cycle
basin of the augmented square-digit function contains a translate of
every finite set of positive integers. In particular there are arbitrarily
long consecutive runs attracted to any specified periodic point, proving
Conjecture 3.3 of Baker Swart et al. The new step constructs an adjacent
pair in every basin using a congruence modulo $(b-1)^2$ and an explicit
base-$b$ carry pattern. A pair-reduction and digit-prefix lifting
argument then gives all finite lengths.
\end{abstract}

\section*{Task 1: the complete consecutive-attraction theorem}

Fix odd integers $b>2$ and $c>0$. Put $B=b-1$, let $\sigma(n)$ denote
the sum of the squares of the base-$b$ digits of $n$, with $\sigma(0)=0$,
and put $S(n)=c+\sigma(n)$ for $n>0$. A cycle $C$ means a periodic orbit
of $S$. Its basin $A_C$ consists of the positive integers whose iterates
eventually enter $C$. In particular every point of $A_C$ eventually
visits every prescribed $u\in C$.

We prove the following strengthening of Conjecture 3.3 of Baker Swart
et al., \emph{Augmented generalized happy functions},
\href{https://arxiv.org/abs/1410.0297}{arXiv:1410.0297v2}; Rocky Mountain
J. Math. 47 (2017), 403--417.
\begin{theorem}
For each cycle $C$ and each finite $T\subset\mathbb Z_{>0}$, there is
$n\ge0$ such that $T+n\subset A_C$. Consequently, for every $u$ on any
cycle, there are arbitrarily long sequences of consecutive $u$-attracted
integers.
\end{theorem}
The even-difference reduction and digit-prefix lifting below reproduce
the mechanisms behind the source's Theorem 3.7 and Lemmas 3.5 and 4.3.
The opposite-parity reduction and the construction of adjacent points
in each basin supply the missing odd-$c$, odd-$b$ case.

\section*{Reduction of a pair}
Let $I(n)=n+1$. An admissible map is a finite composition of $S$ and
nonnegative integer powers of $I$.
Since $b,c$ are odd,
\begin{equation}\label{parity}
S(n)\equiv n+1\pmod2.
\end{equation}
Thus an admissible map changes the parity of every input in the same way.

\begin{lemma}\label{pair}
For distinct positive integers $a<a+w$, an admissible map takes this
pair to two values whose difference in absolute value is
$\varepsilon=w\bmod2\in\{0,1\}$.
\end{lemma}
\begin{proof}
The integer $\sigma(w-1)+1+\varepsilon$ is even. Hence we can choose
$0\le j<B/2$ such that
\[
2j\equiv-\sigma(w-1)-1-\varepsilon\pmod B.
\]
Choose $r\ge1$ so large that
\[
b^r>w-1,\qquad (j+1)b^r-1\ge a,\qquad
rB^2>2j+1+\sigma(w-1).
\]
Translate $a$ to $(j+1)b^r-1$ and then apply $S$. The translated pair
has the digit expansions
\[
j b^r+(b^r-1),\qquad (j+1)b^r+(w-1),
\]
so its image difference, second minus first, is
\[
\Delta=2j+1+\sigma(w-1)-rB^2<0,
\qquad \Delta\equiv-\varepsilon\pmod B.
\]
After ordering the images write them as $q,q+W$, where
$W=-\Delta=Bv+\varepsilon$ with $v\ge0$. Choose $R\ge1$ large enough
that $b^R+v\ge q$ and $b^R>bv+\varepsilon$. Translate this pair to
\[
b^R+v,\qquad b^R+bv+\varepsilon.
\]
There are no carries between the displayed parts, and
$\sigma(bv+\varepsilon)=\sigma(v)+\varepsilon$.
One more application of $S$ therefore leaves difference $\varepsilon$.
Every translation used was nonnegative.
\end{proof}

\begin{corollary}\label{collapse}
An admissible map sends any nonempty finite set to a singleton or a
pair of consecutive integers. A set of one parity can be sent to a
singleton.
\end{corollary}
\begin{proof}
While two distinct members have the same parity, apply Lemma~\ref{pair}
to that pair and apply its map to the whole set. The cardinality drops.
Equation~\eqref{parity} ensures that the only possible terminal sets
have at most one point of each parity. If two remain, apply the odd
case of Lemma~\ref{pair}.
\end{proof}

\section*{Lifting translations back through digit sums}
Fix a cycle $C$. Call a finite set $T$ $C$-good if some nonnegative
translate of $T$ is contained in $A_C$.

\begin{lemma}\label{lift}
If $F$ is admissible and $F(T)$ is $C$-good, then $T$ is $C$-good.
Every singleton is $C$-good. Every finite set of one parity is $C$-good.
\end{lemma}
\begin{proof}
For $F=I$ the assertion follows by increasing the translation by one.
For $F=S$, suppose $S(T)+n\subset A_C$. Choose $r$ with $b^r>\max T$
and let $H$ be the integer whose base-$b$ expansion consists of $n$
ones (take $H=0$ for $n=0$). With $M=Hb^r$ the digit blocks are disjoint,
and
\[
S(t+M)=S(t)+\sigma(H)=S(t)+n\in A_C\quad(t\in T).
\]
Thus $T+M\subset A_C$. Induction proves the assertion for compositions.

For a singleton $\{t\}$ take $z\in C$ and $r\ge0$ with $zb^r\ge t$.
Since $S(zb^r)=S(z)\in C$, translation by $zb^r-t$ works. The final
assertion follows from Corollary~\ref{collapse} and the first assertion.
\end{proof}

\section*{An adjacent pair in every basin}
\begin{lemma}\label{seed}
For each cycle $C$, there are arbitrarily large $N$ such that both
$N$ and $N+1$ belong to $A_C$.
\end{lemma}
\begin{proof}
Apply Lemma~\ref{lift} to a set of $B/2$ consecutive terms of an
arithmetic progression of difference two. We obtain $B/2$ such terms
in $A_C$, covering all residues of one parity modulo $B$. If that
parity is odd, choose a term $x\equiv1\pmod B$; if it is even, choose
$x\equiv2\pmod B$. Put $y=S(x)$. Then $x,y\in A_C$ have opposite
parities, and
\begin{equation}\label{gcd}
\gcd(x,y,B)=1:
\end{equation}
in the first case $\gcd(x,B)=1$, and in the second $\gcd(x,B)=2$
while $y$ is odd. This also covers $B=2$.

Let $w$ be the unique odd integer in $\{1,3,\ldots,B-1\}$ congruent
to $x-y$ modulo $B$, and put $j=(w-1)/2$. By \eqref{gcd} there are
$r_0,s_0\in\{0,\ldots,B-1\}$ solving
\[
xr_0-ys_0\equiv\frac{w-x+y}{B}\pmod B.
\]
For any $r\equiv r_0\pmod B$ and $s\equiv s_0\pmod B$, the binomial
identity $b^t=(1+B)^t\equiv1+tB\pmod{B^2}$ gives
\begin{equation}\label{scaled}
X-Y:=xb^r-yb^s\equiv w\pmod{B^2}.
\end{equation}
First increase $r$ by a multiple of $B$ until
$h=X-c-(j+1)^2$ is positive and as large as desired. Then increase
$s$ by a multiple of $B$ until $Y>X$. This preserves \eqref{scaled}.
Set
\[
k=\frac{Y-X+w}{B^2}\in\mathbb Z_{>0}.
\]
Let $H$ have exactly $h$ base-$b$ digits, all equal to one, and set
\[
N=Hb^{k+1}+jb^k+(b^k-1).
\]
Because $0\le j<b-1$, the displayed digits give exactly
\begin{align*}
S(N)&=c+h+j^2+kB^2=X-w+kB^2=Y,\\
S(N+1)&=c+h+(j+1)^2=X.
\end{align*}
Multiplication by a power of $b$ preserves basin membership, since
$S(tb^m)=S(t)$. Thus $X,Y\in A_C$, whence $N,N+1\in A_C$.
As $h$ can be made arbitrarily large, so can $N$.
\end{proof}

\section*{Completion of the argument}
Given nonempty finite $T$, Corollary~\ref{collapse} supplies an admissible
$F$ such that $F(T)$ is a singleton or a consecutive pair. Singletons
are $C$-good by Lemma~\ref{lift}. A pair $\{a,a+1\}$ is $C$-good by
Lemma~\ref{seed}: choose $N\ge a$ there and translate by $N-a$.
Lemma~\ref{lift} now shows that $T$ is $C$-good. Empty $T$ is trivial.
Taking $T=\{1,\ldots,L\}$ proves the required assertion for every
length $L$ and every $u\in C$. The iterate at which each number visits
$u$ may depend on that number, exactly as the source definition allows.
\qed

\section*{Verification and source scope}
The proof is uniform in $b,c,L$ and independent of computational bounds.
The standard-library verifier checks 20,000 explicit pair reductions,
135 reductions of finite sets, 1,125 digit-prefix identities, and 259
cycle seed constructions across 110 parameter pairs. Of these seeds,
42 are expanded as actual integers; the others are checked using exact
digit-run lengths because their expansions are enormous. There is no
floating-point arithmetic.

All 14 pages of the current author manuscript, arXiv:1410.0297v2
(14 September 2016), were read. A coauthor's institutional CV, created
9 February 2016, identifies this arXiv manuscript as accepted by the
Rocky Mountain Journal of Mathematics, so the current revision is a
post-acceptance author version. The typeset 2017 article could not be
retrieved through the checked public links. These access limits and the
bounded public-claim search are recorded in the accompanying literature
audit. This work has author-agent review, without independent peer review.

\section*{Proposed Tasks 2--5}
\begin{enumerate}
\setcounter{enumi}{1}
\item Derive an explicit computable upper bound, in $b,c,u,L$, for the
least starting point of a length-$L$ consecutive $u$-attracted run by
tracking all translations, exponents and prefix lengths in this proof.
\item For augmented $e$th-power digit sums with $e\ge3$, characterize
the triples $(e,b,c)$ for which every cycle basin contains arbitrarily
long consecutive intervals, including all congruence obstructions.
\item For fixed odd $b,c$, characterize the finite words in cycle-basin
labels that occur on arbitrarily large intervals of consecutive integers.
The present theorem settles constant words only.
\item For each cycle basin with odd $b,c$, determine its lower and upper
natural densities and characterize when those two values coincide.
\end{enumerate}

\begin{thebibliography}{9}
\bibitem{source} B.~Baker Swart, K.~A.~Beck, S.~Crook,
C.~Eubanks-Turner, H.~G.~Grundman, M.~Mei and L.~Zack,
\emph{Augmented generalized happy functions}, Rocky Mountain J. Math.
\textbf{47} (2017), no.~2, 403--417.
\href{https://doi.org/10.1216/RMJ-2017-47-2-403}{DOI:10.1216/RMJ-2017-47-2-403};
\href{https://arxiv.org/abs/1410.0297}{arXiv:1410.0297v2}.
\bibitem{mathdb} \href{https://mathdb.com/p/328940}{MathDB 328940},
\emph{Arbitrarily long consecutive sequences of attracted numbers for
augmented generalized happy functions}. Status checked 28 September 2026.
\end{thebibliography}
\end{document}
